%% ma: L11-B02-0009
%% lop: 11
%% chuong: 1
%% bai: 2
%% dang: 11.2.6
%% loai: TLN
%% muc: TH
%% dap_an: "1,3"
%% nguon: "de-cuong-giua-ky-1-toan-11-nam-2026-2027-truong-thpt-hoang-van-thu-ha-noi-2-8.pdf (M08, Đề cương giữa kì 1 Toán 11 THPT Hoàng Văn Thụ, Hà Nội), Đề minh họa, Phần III Câu 3"
%% kien_thuc: "Bài 1 (giá trị lượng giác của góc lượng giác), Bài 2 (công thức lượng giác)"
%% trang_thai: cho-duyet
%% ly_do: "Dạng mới đề xuất 11.2.6"
\begin{ex}
Cho $\cos2\alpha=-\dfrac{7}{25}$ và $0<\alpha<\dfrac{\pi}{2}$, tính $\tan\alpha$ (kết quả làm tròn đến hàng thập phân thứ nhất).
\shortans{1,3}
\loigiai{$\sin^2\alpha=\dfrac{1-\cos2\alpha}{2}=\dfrac{16}{25}$, $\cos^2\alpha=\dfrac{9}{25}$. Do $0<\alpha<\dfrac\pi2$: $\tan\alpha=\dfrac43\approx1{,}3$.}
\end{ex}
